\documentclass[11pt,a4paper]{article}
% Shared preamble for RuPhO English problem statements (match T1 2026 style).
% Use: \input{latex/rupho_en_preamble.tex} after \documentclass.
\usepackage[margin=1in]{geometry}
\usepackage{amsmath,amssymb}
\usepackage{graphicx}
\usepackage{float}
\usepackage{enumitem}
\usepackage{microtype}
\usepackage{caption}
\usepackage{tabularx}
\usepackage{hyperref}

\captionsetup{font=small,labelfont=bf,skip=6pt}

\setlist[itemize]{leftmargin=1.2cm}
\setlist[enumerate]{leftmargin=1.2cm}

% One outer box per subproblem: [id | statement | points] (no inner rules)
\newcommand{\problembox}[3]{%
  \par\addvspace{0.42em}%
  \noindent
  \setlength{\fboxsep}{7.5pt}%
  \setlength{\fboxrule}{0.95pt}%
  \fbox{%
    \begin{minipage}{\dimexpr\linewidth-2\fboxsep-2\fboxrule\relax}
    \begin{tabularx}{\linewidth}{@{}>{\raggedright\arraybackslash\bfseries}p{2.1em} @{\hspace{0.9em}} >{\raggedright\arraybackslash}X @{\hspace{0.9em}} >{\raggedleft\arraybackslash\bfseries}p{2.4em}@{}}
    #1 & #2 & #3
    \end{tabularx}
    \end{minipage}%
  }%
  \par\addvspace{0.32em}%
}
\title{\textbf{Halbach magnetic assembly --- Solution}}
\author{\normalsize X20 (2020) --- English translation}
\date{}
\begin{document}
\maketitle
\problembox{A1}{The magnetic field weakens with increasing distance; for a given angle $\theta$, find the dependence of the magnetic field $B$ at a distance $r$ from the dipole.\\ The dependence of the magnetic field of a dipole on $\vec r$: $\vec B (\vec r) = \frac {\mu_0}{4\pi} (\frac{3\vec r (\vec m \cdot \vec r )} {r^5} - \frac { \vec m} {r^3}), \vec m \cdot \vec r = mr \cos \theta \Rightarrow$$\vec B (\vec r) = \frac {m \mu_0}{4\pi r^3}(3\cos \theta \hat{r} - \hat{y}) $ The magnitude of this vector can be found in different ways, for example, it is convenient to introduce an additional unit vector: $\hat{\theta} = - \sin \theta \cdot \hat{y} + \cos \theta \cdot \hat{x}, \hat{r}=\cos \theta \cdot \hat{y} +\sin \theta \cdot \hat{x}$ Or otherwise: $\hat{y}=\cos\theta\cdot \hat{r} - \sin \theta \cdot \hat{\theta}, \hat{x}= \sin \theta \cdot \hat{r} + \cos \theta \cdot \hat{\theta}$\\ $3 \cos \theta~\hat{r} - \hat{y} = 3 \cos \theta~\hat{r} - (\cos \theta \cdot \hat{r} - \sin \theta \cdot \hat{\theta}) \Rightarrow$$\vec B (\vec r) = \frac {\mu_0 m} {4 \pi r^3} ((2 \cos \theta) \hat{r} + (\sin \theta) \hat{\theta}) \Rightarrow B = \frac {\mu_0 m}{4\pi r^3} \sqrt{(2\cos \theta)^2+(\sin \theta)^2)} $}{0.50}
\problembox{B1}{Express the magnetic field $B(y)$ along an axis perpendicular to the magnet at a distance $y$ from the center.\\ To find the answer to this point, we divide the magnet into infinitely narrow rings and integrate:\\ We obtain the following expression: $\vec B (y) = \frac {\mu_0} {4\pi} \cdot \int ^R_0 2\pi x \cdot dx \cdot \sigma \cdot \frac {1} {(x^2+y^2)^{\frac{3}{2}}} \cdot (2 \cos \theta \cdot \cos \theta \cdot \hat{y} + \sin \theta \cdot \sin \theta \cdot (-\hat{y}))$ For convenience, we proceed to integration over the angle: $R = y \cdot \tan \theta_{max}, x=y\cdot \tan \theta, dx=y\cdot \frac {d\theta}{\cos^2\theta}, \frac{y}{\cos\theta}=\sqrt{x^2+y^2}$ Then $B=\frac {\mu_0 \sigma}{2} \cdot \int ^{\theta_{max}}_0 y\cdot \tan \theta \cdot y \cdot \frac {d\theta} {\cos^2\theta} \cdot \frac{\cos^3 \theta}{y^3}\cdot \hat{y} \cdot (2\cos^2\theta-\sin^2\theta)$$=\frac{\mu_0\sigma \cdot \hat{y}}{2y} \cdot \int^{\theta_{max}}_0 d\theta \cdot (2 \cos^2\theta \sin \theta - \sin^3 \theta)$$=\frac{\mu_0\sigma \cdot \hat{y}}{2y} \cdot \int^{\theta_{max}}_0 d \theta \cdot (3\cos^2 \theta \sin \theta - \sin \theta)$$=\frac{\mu_0\sigma\cdot \hat{y}}{2y} \cdot [-\cos^3\theta+\cos\theta]|^{ \theta_{max}}_0=\frac{\mu_0\sigma \cdot \hat{y}}{2y}\cdot[\cos \theta_{max}- \cos^3 \theta_{max}]$$\cos \theta_{max} = \frac {y} {\sqrt{R^2+y^2}}$ Simplifying, we obtain the final result: $B(y)=\frac{\mu_0\sigma}{2}\cdot \frac{R^2}{(R^2+y^2)^{\frac{3}{2}}}$}{1.50}
\problembox{B2}{Estimate the magnitude of the magnetic field near the surface of the magnet. Express your answer in terms of quantities $t,~D,~\rho,~\mu_0$.\\ Since the relation $t \ll D$ is satisfied, the magnet can be considered flat and its thickness can be neglected when the field is found at the surface: $\sigma=\rho\cdot t, B(y)=\frac{ \mu_0\sigma}{2}\cdot \frac {R^2} {(R^2+y^2)^{\frac{3}{2}}}, y=0 \Rightarrow B_0 = \frac {\mu_0\rho t} {D}$$\frac{\mu_0\rho t}{D}=0.13 \text{T}$}{0.50}
\problembox{B3}{Obtain the expression and numerical value of the interaction force $F_0$ between the door and the magnet pressed to it, and also calculate the pressure $P_0$ of the magnet on the door.\\ B3) According to the condition, the volumetric energy density of the magnetic field is: $\frac{E}{V}=\frac {1}{2\mu_0} \cdot B^2$ When the magnet is separated from the door a short distance $y$: $\Delta E = \pi (D/2)^2\cdot y\cdot \frac{1}{2\mu_0} \cdot B^2_0 $$F=\frac {\Delta E}{y} = \pi (D/2)^2 \cdot \frac {1}{2\mu_0}\cdot B^2_0= 2.2 H$ Knowing the force and area of contact, it is easy to express the pressure: $P=\frac{F}{S}=\frac{B^2_0}{2\mu_0} =6.9\cdot 10^3~\text{Pa}$}{0.50}
\problembox{C1}{Write down an expression for the field $\vec B(\vec r_0, y)$ that is created by a series of magnets. (For convenience, the field is expressed as both $\vec r_0$ and $y$, although technically it is $y = (\vec r_0)_y$.)\\ According to the diagram below:\\ $r_0=\frac{y}{\cos \alpha}, x=y\tan \alpha, \vec{r_0} = y\hat{y} + x\hat{x} \Rightarrow r_0 = \hat{y}\cdot \cos \alpha +\hat{x} \cdot \sin \alpha$$\frac{r_0}{r}=\cos \theta, z=r_0 \tan \theta \Rightarrow dz =\frac{r_0d\theta} {cos^2\theta}$$\vec r=x\hat{x}+y\hat{y}+(-z)\hat{z}$$\vec B(\vec{r_0},y)=\frac{\mu_0}{4 \pi}\cdot \int^\infty_{-\infty}dz \cdot \rho_L\cdot[\frac{3}{r^4} (\frac{y\cdot \cos \theta}{r_0}) \cdot (\vec{r_0}-z\hat{z})-\frac{\hat{y}}{r^3}]= $$\frac{\mu_0\cdot\rho_L}{4\pi}\cdot \int^{\pi/2}_{-\pi/2}d\theta \cdot \frac{r_0}{\cos^2\theta}\cdot[ \frac{3y}{r_0^5}\cdot \cos^5\theta \cdot \vec{r_0} - \frac {\hat{y}}{r_0^3} \cdot \cos^3\theta]=$$\frac{\mu_0\cdot\rho_L}{4\pi r_0^2}\cdot \int^{\pi/2}_{-\pi/2}d\theta\cdot[\frac{3y\cdot\vec{r_0}}{r_0^2} \cdot\cos^3\theta-\hat{y}\cdot\cos\theta]=$$\frac{\mu_0\cdot\rho_L}{4\pi r_0^2}\cdot [\frac{3y\cdot\vec{r_0}}{r_0^2} \cdot(\sin\theta-\frac{1}{3}\sin^3\theta)- \hat{y}\cdot\sin\theta]|_{-\pi/2}^{\pi/2}=$$\frac {\mu_0\cdot\rho_L} {4\pi r_0^2}\cdot[\frac{4y\cdot\vec{r_0}}{r_0^2}-2\hat{y}]$}{2.00}
\problembox{C2}{Find the magnetic field on both sides of the assembly. Give the answer in the form of some integral.\\ According to the diagram below:\\ $r_0=\frac{y_0}{\cos \alpha}=\frac{y}{\cos( \beta+\alpha)} \Rightarrow  y=y_0\cdot \frac {\cos(\beta+\alpha)}{\cos\alpha}$$\hat{y}=\cos \beta \cdot \hat{y_0} +\sin \beta \cdot \hat{x}, \hat{r_0}=\cos \alpha \cdot \hat{y_0} - \sin \alpha \cdot \hat{x}$$\vec B = \int_{-\infty}^{\infty} dx \cdot \frac {\mu_0\sigma}{4\pi y_0^2} \cdot \cos^2 \alpha \cdot [4\cos(\beta+\alpha)[\cos\alpha \cdot \hat{y_0}-\sin\alpha\cdot\hat{x}]- 2[\cos \beta \cdot \hat{y_0}+\sin \beta \cdot \hat{x}] ]$ Through simple mathematical transformations we get: $\frac{\mu_0\sigma}{2\pi y_0}\cdot \int _{-\pi/2}^{\pi/2} d\alpha\cdot [\hat{y_0}\cdot \cos(\beta+2\alpha)-\hat{x}\cdot\sin(\beta+2 \alpha)]$ Substitute the dependence for $\beta$: $\beta=\beta_0+kx_0+ky_0\cdot\tan \alpha$ We get: $\frac{\mu_0\sigma}{2\pi y_0}\cdot \int_{-\pi/2}^{\pi/2}d\alpha\cdot [\hat{y_0}\cdot\cos(\beta_0+kx_0+ky_0\cdot\tan\alpha +2\alpha)-\hat{x}\cdot\sin(\beta_0+kx_0+ky_0\cdot\tan\alpha+2\alpha)]$}{1.00}
\problembox{C3}{Show that on one side of an ideal assembly the magnetic field tends to zero.\\ Let's look carefully at the resulting expression: $\cos(\beta_0+kx_0+ky_0\cdot\tan\alpha+2\alpha)$$=\cos(ky_0\cdot\tan\alpha+2\alpha)-\sin(\beta_0+kx_0) \sin(ky_0\cdot\tan\alpha+2\alpha)$$\cos(ky_0\cdot\tan\alpha+2\alpha)=\cos(ky_0\cdot\tan \alpha)\cos(2\alpha)-\sin(ky_0\cdot\tan\alpha)\sin (2\alpha)$ By condition, $\int_{-\pi/2}^{\pi/2} dx\cdot\cos(2x)\cos(c\cdot \tan x)=\frac{c\cdot\pi}{e^c}= \int^{\pi/2}_{-\pi/2}dx\cdot\sin(2x)\cdot \sin(c\cdot\tan x)$ It is easy to see that this is our integral, if we take the value $ky_0$ as $c$.}{1.00}
\problembox{C4}{Write down the expression for the field on the other side.\\ When moving to the other side of the assembly, some signs in the equation change to the opposite: $\vec{B}=-\frac{\mu_0\sigma}{2\pi y_0} \cdot \int^{\pi/2}_{-\pi/2} d\alpha \cdot [\hat{y_0}\cdot [\cos(\beta_0+kx_0)\cdot \cos(-ky_0\cdot\tan\alpha+2\alpha)-\sin"..."sin"..."] -\hat{x}\cdot[\sin(\beta_0+kx_0)\cos(-ky_0\cdot \tan\alpha+2\alpha)+\cos"..."sin"..."] ]$ Integrals from odd functions will vanish due to symmetry considerations, even functions remain. $\vec B=- \frac{\mu_0\sigma}{2\pi y_0} \cdot \int^{\pi/2} _{-\pi/2}d\alpha\cdot[\hat{y_0}\cdot[\cos(\beta_0+kx_0-ky_0 \cdot\tan\alpha+2\alpha)]-\hat{x}\cdot[\sin(\beta_0+kx_0) \cos(-ky_0\cdot\tan\alpha+2\alpha)+\cos"..."\sin"..."] ]$$=-\frac{\mu_0\sigma}{2\pi y_0}\cdot\int ^{\pi/2}_{-\pi/2}d\alpha\cdot[\hat{y_0}\cdot \cos(\beta_0+kx_0)\cdot(cos(2\alpha)\cos(ky_0 \cdot\tan\alpha)+\sin"..."\sin"...")]- \hat{x}\cdot[\sin(\beta_0+kx_0)\cdot"..."] ]$$=-\frac{\mu_0\sigma}{2\pi y_0}\cdot [\hat{y_0}\cdot\cos(\beta_0+kx_0)-\hat{x}\cdot \sin(\beta_0+kx_0)]\cdot 2 \cdot\frac{ky_0\pi} {e^{ky_0}}$$=-\mu_0\sigma k\cdot e^{-ky_0}\cdot[\hat{y_0} \cdot\cos(\beta_0+kx_0)-\hat{x}\cdot \sin(\beta_0+kx_0)]$}{1.00}
\problembox{C5}{Based on the field expression, find the average pressure $$ of such a magnet on the refrigerator door. Take the following parameters: thickness $t=0.5~\text{mm}$, volumetric density of the magnetic dipole $\rho = 2 \cdot 10^5 \frac {\text{T} \cdot \text{m}} {\text{H}}$, assembly step $\lambda=5~\text{mm}$.\\ The force and pressure are the same as in point B3: $y_0 \rightarrow 0 \Rightarrow \vec B = -\mu_0\sigma k \cdot[\hat{y_0}\cdot\cos(\beta_0+kx_0) -\hat{x}\cdot\sin(\beta_0+kx_0)] \Rightarrow$$B^2=(\mu_0\sigma k)^2, P=\frac{1}{2\mu_0}\cdot B^2, \sigma=\rho t \Rightarrow$$B=(1.257\cdot 10^{-6}~\text{H/m}) \cdot(2\cdot 10^5~ \text{T}\cdot\text{m}/\text{H}) \cdot(5\cdot10^{-4}~\text{m}) \cdot \frac{2\cdot 3.14}{5\cdot 10^{-3}\text{ m}} =0.16 \text{ T}$$P=\frac{1}{2\mu_0}\cdot B^2=\frac{1} {2\cdot 1.257 \cdot 10^{-6} \text{ H/m}} \cdot (0.16\text{ T}^2)=10\text{ kPa}$}{1.50}
\problembox{C6}{Find the ratio between the pressure created by a Halbach magnetic assembly and the pressure created by a conventional magnet of the same material, with the same radius and thickness. Here, too, we should neglect the effects on the perimeter of the circle and the thickness of the magnet.\\ According to the formulas obtained earlier: $\eta=\frac{P}{P_0}=(\frac{B}{B_0})^2 (\frac{\mu_0\rho tk}{\mu_0\rho t/2R})^2= (\frac{4\pi R}{\lambda})^2$}{0.50}
\vspace{1em}
\noindent\textit{Source:} \url{https://pho.rs/p/19/s}
\end{document}